\section{Learning the Value of a Released Edit}

Timing is not equivalent to changing the recipients.
We fix the CV--OLS disagreement order at each query and let
each action choose a prefix. A two-bin FIFO state records old
credits; the current release is observed, but the end of the
recording is not. On donors, backward dynamic programming
labels the advantage of each spend count over spending zero,
using an additive lower-gain proxy with positive risk costs and
a motion-scaled penalty for expiry. This objective is not a
certificate for the original recording-level constraints.
The same head also predicts the immediate advantage without
the continuation term. Candidate and release heads used on
donors exclude both the fitting row's site and the outer site.

\begin{table}[h]
\centering
\caption{Learned timing on the same frozen neural bank with
a fixed causal recipient ranking. Joint passes use the
unchanged recording-level risk and both lower gains.}
\small
\begin{tabular}{lrrrr}
\toprule
Policy & Edits & Unused & Undefined & Joint \\
\midrule
Immediate & \QueueValueNowEdits{} & 0 & \QueueValueNowUndefined{} & \QueueValueNowJoint{} \\
Myopic & \QueueValueOneEdits{} & \QueueValueOneUnused{} & \QueueValueOneUndefined{} & \QueueValueOneJoint{} \\
Continuation & \QueueValueTailEdits{} & \QueueValueTailUnused{} & \QueueValueTailUndefined{} & \QueueValueTailJoint{} \\
\bottomrule
\end{tabular}
\end{table}

The continuation policy yields \QueueValueDifference{} versus
the myopic arm and \QueueValueVersusCV{} versus CV.
Against immediate use of the same realized credits, its effect
is \QueueValueContinue{}; the myopic arm yields \QueueValueMyopic{}.
Intervals resample sites after averaging seeds, not windows.
Neither policy establishes a reliable gain or passes the
full recording roster.

To separate the action menu from learning, we also solve
the exact label-aware problem under this fixed ranking:
at most one prefix is chosen per query, its cardinality
enters the original cumulative budget, and all credits
must be spent. The original risk, support and dual-gain
conditions hold for \QueueValueFixed{} (zero delay) and
\QueueValueDelay{} (two queries) of \QueueValueViews{} views.
Removing the fixed ranking raises the two-query ceiling to
\NeuralCreditShort{}. The loss is attributable to the restricted
action menu under this budget, not solely to timing estimation.
Timing and recipient decisions remain coupled.
